\section{The Data Kernel Perspective Space}
For our purposes, a model $ f \in \mathcal{F} $ is a random mapping from a query space $ \mathcal{Q} $ to a response space $ \mathcal{X} $. 
Given $ q \in \mathcal{Q} $, model responses $ f(q)_{1}, \hdots, f(q)_{r} $ are sampled $ i.i.d. $ from the distribution $ F $.
We let $ g: \mathcal{X} \to \mathbb{R}^{p} $ be a fixed embedding function that maps a response to a real-valued vector.

Given models $ f_{1}, \hdots, f_{n} $ and queries $ Q = \{q_{1}, \hdots, q_{m}\} $,  we let $ \bar{X}_{i} \in \mathbb{R}^{m \times p} $ be the matrix whose $ j $th row is the average embedded response from $ f_{i} $ to query $ q_{j} $; or,  $\bar{X}_{ij\cdot} = \frac{1}{r} \sum_{k=1}^{r} g(f_{i}(q_{j}))_{k} $.
Further, we let $ D $ be the pairwise distance matrix with entries $ D_{ii'} = || \bar{X}_{i} - \bar{X}_{i'} ||_{F} $.
We refer to the distribution on embedded responses induced by $ f_{i}(q_{j}) $ as $ F_{ij} $.

Following \citep{aranyak}, the $d$-dimensional Data Kernel Perspective Space (DKPS) representations of the models are defined as the vectors $ (\widehat{\psi}_{1}, \hdots, \widehat{\psi}_{n}) $ that are a solution to
\begin{equation}
(\widehat{\psi}_{1}, \hdots, \widehat{\psi}_{n}) =    \text{argmin}_{z_i \in \mathbb{R}^d} 
    \sum_{i,i'}^{n}
    (
\lvert\lvert
z_i - z_{j} \rvert\rvert
 -
D_{ii'}
    )^2.
\label{eq:dkps}
\end{equation}
DKPS enables treating model-level analysis as a problem in a low-dimensional Euclidean space.
We note that the solution to Eq. \eqref{eq:dkps} -- and the quality of the representations for a given task -- may depend on the choice of query set $ Q $ and embedding function $ g $.
When the context is not clear, we emphasize the dependence on $ Q $ by referring to $ \psi $ as $ \psi(Q) $.
In the benchmark setting, the choice of $ Q $ is reasonably constrained to elements of $ 2^{Q^{*}} $.
Finally, we let $ \widehat{\Psi}_{Q}: \mathcal{F} \to \mathbb{R}^{d} $ be a mapping from the model space to the estimated perspective space under $ Q $; that is $ \widehat{\Psi}_{Q}(f_{i}) = \widehat{\psi}_{i}(Q) $.

Figure \ref{fig:dkps-example} shows the $ d = 2 $ DKPS of models induced by various choices of $ n $ and $ m $. 
Each dot is a model colored by score on the counting and probability subtask from HELM-Lite.
\begin{figure}
    \centering
    \includegraphics[width=\linewidth]{figures/fig1.pdf}
    \caption{Example $ d = 2 $-dimensional Data Kernel Perspective Spaces (DKPS) for models publicly evaluated on HELM-Lite's MATH counting and probability subtask. 
    Each panel includes the DKPS representations for different $(n,m)$ =
     (number of models, number of queries) pairs induced by a random query set of size $ m $.
    Each dot is a model colored by its score on the subtask.
    As the number of queries increases (left to right), the models with similar scores are more tightly clustered.
    As the number of models increases (top to bottom), there are more models to help localize score signal.}
\label{fig:dkps-example}
\end{figure}

\subsection{Theoretical properties of the DKPS}
For a fixed collection of models, query set, and number of responses, the $ d $-dimensional DKPS representations of the models -- $ \widehat{\psi}_{1}, \hdots, \widehat{\psi}_{n} $ --  are estimates of a set of ``true" $ d $-dimensional vectors $ \psi_{1}, \hdots, \psi_{n} $.
That is, as $ n, m, $ and $ r $ tend to infinity at particular relative rates, $ \lim \widehat{\psi}_{i} \to \psi_{i} $ \citep{aranyak}.
Further, the error of the worst estimate is bounded above by a constant dependent on the number of models, the number of queries, the number of replicates, and the DKPS dimension: $ \max_{i}|| \widehat{\psi}_{i} - \psi_{i}||_{2} \le c(n,m,r,d)$ with high probability \citep{acharyya2025concentrationboundsresponsebasedvector}.
Given its importance to provable query-efficiency, we state this last result in its entirety here:
\begin{theorem}
\label{thm:concentration}
  [\citet{acharyya2025concentrationboundsresponsebasedvector} Theorem 2] In our setting, suppose $ r = \omega(n^{3}) $ and $ \sup_{ij} \;\text{trace}\left(Cov\left(F_{ij}\right)\right) = O(1)  $. 
  Then, under technical assumptions, for every $ \delta \in (0, 1/2) $ and for sufficiently large $ n $ and $ r $  \begin{align*}
      \max_{i} || \widehat{\psi}_{i}  - \psi_{i} ||_{2} \le \text{Poly}_{3}\left(\left(\frac{n^{3}}{r}\right)^{\frac{1}{2} - \delta}\right)
  \end{align*}
  with high probability, where $ \text{Poly}_{3} \left(x\right) $ is a third degree polynomial in $ x $.
\end{theorem}
The concentration result of Theorem \ref{thm:concentration} is necessary to make claims related to the quality of representations for a given $ (n, m, r) $, which is critical for analyzing the query efficiency of DKPS-based methods for benchmark score prediction.

\subsection{Inference in the DKPS}
Following \citet{helm-etal-2025-statistical}, we consider model-level inference in DKPS as a statistical problem: Let $ (f, y), (f_{1}, y_{1}), \hdots, (f_{n}, y_{n}) $ be $ i.i.d. $ samples from a joint distribution on (model,  score) pairs. 
In the benchmark prediction setting, given $ f $ the true benchmark score $ y $ is deterministic and analysis with the marginal distribution on $ f $ is equivalent to analysis with the joint distribution. 
We refer to the distribution on models as $ P_{f} $.

Our goal is to construct a decision function $ h: \mathcal{F} \to \mathcal{Y} $  that minimizes the expected loss over $ P_{f} $. 
In practice, operating directly on the space of models is too complex.
Instead, we consider the proxy problem of model-level inference in DKPS: we observe perspective-score pairs $ (\widehat{\Psi}_{Q}(f), y), (\widehat{\Psi}_{Q}(f_{1}), y_{1}), \ldots, (\widehat{\Psi}_{Q}(f_{n}), y_{n}) $ and use these to construct a decision function $ h_{n}^{Q}: \mathcal{F} \to \mathcal{Y} $ defined by $ h_{n}^{Q}(f) = \hat{h}_{n}^{Q}(\widehat{\Psi}_{Q}(f)) $ where $ \hat{h}_{n}^{Q}: \mathbb{R}^{d} \to \mathcal{Y} $ is trained on $ \{(\widehat{\Psi}_{Q}(f_{i}), y_{i})\}_{i=1}^{n} $.
Formally, with loss function $ \ell: \mathcal{Y} \times \mathcal{Y} \to \mathbb{R} $, our goal is to select $ \hat{h}_{n}^{Q} $ to minimize
\begin{equation*}
\mathbb{E}_{P_{f}}[\ell(h_{n}^{Q}(f), y)] = \mathbb{E}_{P_{f}}[\ell(\hat{h}_{n}^{Q}(\widehat{\Psi}_{Q}(f)), y)].
\end{equation*}
We sometimes refer to $ h_{n}^{Q} $ as $ h_{n}^{(m)} $ as context requires.

% Our goal is to construct a decision function $ h: \mathcal{F} \to \mathcal{Y} $  that minimizes the expected loss over $ P_{f} $. 
% In practice, operating directly on the space of models is too complex.
% Instead, we consider the proxy problem of model-level inference in DKPS: we observe  $ (\widehat{\Psi}_{Q}(f), y), (\widehat{\Psi}_{Q}(f)_{1}, y_{1}), \hdots, (\widehat{\Psi}_{Q}(f)_{n}, y_{n}) $ and want to construct a decision function $ h_{n} := h(\;\cdot\;; \{(\widehat{\Psi}_{Q}(f)_{i}, y_{i})\}_{i=1}^{n}\}) $ that minimizes the expected loss over the marginal distribution $ P_{f} $. 
% Or, with loss function $ \ell: \mathcal{Y} \times \mathcal{Y} \to \mathbb{R} $, select $ h $ such that $ \mathbb{E}_{\widehat{\psi}(Q)}[\ell(h(\widehat{\psi}(Q)), y)]$ is minimized.

% Given its dependence on $ |Q| = m $ (and more generally $ Q $ itself) we sometimes refer to $ h_{n} $ as $ h_{n}^{Q} $ or $ h_{n}^{(m)} $ as context requires.

\section{Provable query-efficiency in the DKPS}
\label{sec:theory}
In this section we define query-efficiency and prove that a simple regression function on the perspectives is query-efficient relative to estimating the model's score with its score on a subset of the benchmark queries.
Recall $ y := y(f, Q^{*}) $.

For a fixed $ Q \subseteq Q^{*} $ with $ |Q | = m$, we say the sequence of decision functions $ (h_{1}^{Q}, h_{2}^{Q}, \hdots) $ is \textit{$Q $ query-efficient} relative to $ (h_{1}'^{Q}, h_{2}'^{Q}, \hdots) $ if there exists $ N \in \mathbb{N} $ such that
\begin{equation}
\label{eq:query-efficiency}
\small{
\mathbb{E}_{P_f}\,
% \mathbb{E}_{\widehat Q}
\Big[\ell\big(h^{Q}_{n}(f),\; y\big)\Big]
\le\;
\mathbb{E}_{P_{f}}\,
% \mathbb{E}_{\widehat Q}
\Big[\ell\big(h'^{Q}_{n}(f),\; y\big)\Big]
}
\end{equation}
for all $ n > N $. 
We say the sequence of decision functions $ (h_{1}^{(m)}, \hdots, h_{n}^{(m)}) $ is \textit{m query-efficient} relative to $ (h_{1}'^{(m)}, \hdots, h_{n}'^{(m)}) $ if for all $ Q $ with $ |Q | = m $ there exists $ N_{Q} \in \mathbb{N} $ such that Eq. \eqref{eq:query-efficiency} holds for all $ n > N_{Q} $.

Finally, we say the sequence of decision functions $ (h_{1}^{(m)}, \hdots, h_{n}^{(m)}) $ is \textit{query-efficient} relative to $ (h_{1}'^{(m)}, \hdots, h_{n}'^{(m)}) $ if for all $ m < M $ there exists $ N^{(m)} \in \mathbb{N} $ such that it is $ m $-query efficient.


For our theoretical analysis we assume the benchmark score function is well-defined on all subsets of $ Q^{*} $.
For $ Q \subset \mathcal{Q} $, we let $ \hat{y}_{Q} := y(f, Q) $.

Let $ h_{n}^{(m)} $ be nearest neighbor regression in perspective space and $ \delta^{*} = \min_{i} || \widehat{\psi}_{i} - \widehat{\psi}||_{2} $:
\small{
\begin{align*}
\hat{y}_{NN} := h^{(m)}_{n}(\widehat{\psi})
= \frac{\sum_{i=1}^n \mathbf{1}\{i: ||\widehat{\psi}_{i} - \widehat{\psi}||_{2} = \delta^{*} \}y_{i}}{\sum_{i=1}^n \mathbf{1}\{i: ||\widehat{\psi}_{i} - \widehat{\psi}||_{2} = \delta^{*}\}}.
\end{align*}
}

To establish query-efficiency of nearest neighbor regression in DKPS, we require two key assumptions.
\begin{assumption}[Lipschitz Score Function]
\label{ass:lipschitz-score}
Given $ Q \subseteq 2^{Q^*} $, the benchmark score function $y(\;\cdot\;, Q^{*}): \mathcal{F} \rightarrow \mathbb{R}$ is $\gamma$-Lipschitz on the perspective space induced by $ Q $; or, there exists $ \gamma > 0 $ such that for any $f, f' \in \mathcal{F}$,
$$|y(f, Q^{*}) - y(f', Q^{*})| \leq \gamma \cdot ||\psi(Q) - \psi'(Q)||_2.$$
\end{assumption}
The smoothness on the mapping from $ y $ to $ \psi(Q) $ ensures that nearby models in DKPS have similar benchmark scores.
In practice, when $ m = M $ Assumption \ref{ass:lipschitz-score} will hold if the embedding function $ g $ is sufficiently smooth with respect to $ y $.
With the same condition on $ g $,
Assumption \ref{ass:lipschitz-score} will hold in practice for $ m < M $ in a probabilistic sense.

\begin{assumption}[Model Distribution Support]
\label{ass:model-support}
The model distribution $P_f$ has non-zero measure on all compact subsets of $\mathcal{F}$. Equivalently, for any target model $f$ and radius $\delta > 0$, there exists $\epsilon > 0$ such that 
$$P_f(B_\delta(f)) \geq \epsilon$$
where $B_\delta(f) = \{f' \in \mathcal{F}: d_{\mathcal{F}}(f, f') < \delta\}$ for some appropriately defined $ d_{\mathcal{F}} $.
\end{assumption}
Denseness on the space of models ensures that as $n$ increases, we observe models close to any target model with high probability.
In practice, Assumption \ref{ass:model-support} suggests that a large, diverse set of reference models may be needed to realize query efficiency for an arbitrary target model.

Given Assumptions \ref{ass:lipschitz-score} \& \ref{ass:model-support} we are able to arbitrarily bound the prediction error of nearest neighbor regression as a function of $(n, m, r)$.
We now state our main theoretical result:
\begin{theorem}
\label{thm:query-efficiency}
    For any $ \epsilon > 0 $ there exists $ (n,m,r) $ such that 
    %$ MSE(y_{NN}) \le 2\cdot \epsilon^{2} $ 
    $ MSE(\hat{y}_{NN}) \le \epsilon $
    with high probability.
    That is, for $ m < M $ such that $ MSE(\hat{y}_{Q}) > 0 $, $ \hat{y}_{NN} $ is query-efficient relative to $ \hat{y}_{Q} $ with high probability. 
\end{theorem}
We provide the proof of Theorem \ref{thm:query-efficiency} in its entirety.
\begin{proof}[Proof.]
Fix $ Q \subseteq Q^{*} $. 
Let $f$ be a target model with true perspective $\psi$ and estimated perspective $\widehat{\psi}$, and let $f^* \in \arg\min_{i} \|\widehat{\psi}_i - \widehat{\psi}\|_2$ denote one if its nearest neighbors with perspectives $\psi^*$ and $\widehat{\psi}^*$. The prediction error is $|\hat{y}_{NN} - y| = |y(f^*, Q^*) - y(f, Q^*)|$.

By Assumption \ref{ass:lipschitz-score}, we have $|y(f^*, Q^*) - y(f, Q^*)| \leq \gamma \cdot \|\psi^* - \psi\|_2$. 
By the triangle inequality:
\begin{equation*}
\|\psi^* - \psi\|_2 \leq \|\psi^* - \widehat{\psi}^*\|_2 + \|\widehat{\psi}^* - \widehat{\psi}\|_2 + \|\widehat{\psi} - \psi\|_2.
\end{equation*}
Let $\delta^* = \|\widehat{\psi}^* - \widehat{\psi}\|_2$ and let $c = c(n,m,r,d)$ be the concentration bound from Theorem \ref{thm:concentration}. With high probability, $\max_i \|\widehat{\psi}_i - \psi_i\|_2 \leq c$, so $\|\psi^* - \psi\|_2 \leq 2c + \delta^*$. Thus $|y^* - y| \leq \gamma(2c + \delta^*)$.

By Theorem \ref{thm:concentration}, for $r = \omega(n^3)$ and sufficiently large $n $ and $ r$, the constant $c \leq \text{Poly}_3(({n^3}/{r})^{{1}/{2}-\delta})$ can be made arbitrarily small. 

By Assumption \ref{ass:model-support}, for any $\epsilon' > 0$, there exists $n_0$ such that for $n > n_0$, $P_f(\min_i \|\psi_i - \psi\|_2 < \epsilon') \geq 1 - \eta$ for arbitrarily small $\eta > 0$. Combined with the concentration result, $\delta^* < \epsilon' + 2c$ with high probability.

Therefore, with high probability:
\begin{equation*}
|y^* - y| \leq \gamma(2c + \delta^*) < \gamma(4c + \epsilon').
\end{equation*}
For any $\epsilon > 0$, choose $c < \epsilon^{1/2}/(8\gamma)$ and $\epsilon' = \epsilon^{1/2}/(2\gamma)$ with $ n $ and $ r $ sufficiently large.
Then with high probability, $|y^* - y| < \gamma \cdot \epsilon^{1/2}/\gamma = \epsilon^{1/2}$, so $MSE(\hat{y}_{NN}) = \mathbb{E}_{P_f}[(y^* - y)^2] \leq \epsilon$ with high probability.

Thus, if $MSE(\hat{y}_{Q}) > 0$ then there exists $N$ such that for $n > N$, $MSE(\hat{y}_{NN}) < MSE(\hat{y}_{Q})$, establishing query-efficiency.
\end{proof}

That is, benchmark score prediction using the geometry induced by DKPS using a sufficiently large number of models and cached responses to a subset of benchmark queries is better than using a model’s score on the subset of queries.